\documentclass[10pt,journal,compsoc,onecolumn]{IEEEtran}
\usepackage[utf8]{inputenc}
\usepackage{geometry}
\geometry{letterpaper, margin=0.75in}
\usepackage{amsmath}
\usepackage{amsfonts}
\usepackage{booktabs}
\usepackage{array}
\newcolumntype{L}[1]{>{\raggedright\arraybackslash}p{#1}}
\usepackage{float}
\usepackage{graphicx}
\usepackage{tikz}
\usepackage[american]{circuitikz}
\usetikzlibrary{shapes.geometric, arrows, positioning, calc}
\usepackage[hidelinks]{hyperref}
\usepackage{import}
\setcounter{tocdepth}{2}

\begin{document}

\title{Current Line-Following Robot: Construction, Wiring and Control}
\author{Jan Balewski, NERSC, LBNL\\[0.5em]Updated October 11, 2026}
\maketitle

\noindent This document describes the current three-sensor, two-wheel robot and its working local line follower. The ASUS network is configured and SSH is verified. The remote Desktop controller is implemented and deployed; motor-disabled checks and short moving line-following runs are recorded. Clock-delay measurements, the telemetry video and the earlier timing results are in the appendices, with long-term neural workloads, motion simulation and historical prototypes.

\tableofcontents
\subimport{line_folower_figs/}{section_physical_robot.tex}
\subimport{line_folower_figs/}{section_robot_network.tex}
\subimport{line_folower_figs/}{section_electrical_robot.tex}
\subimport{line_folower_figs/}{section_local_line_follower.tex}
\subimport{line_folower_figs/}{section_remote_control.tex}

\clearpage
\appendices
\subimport{line_folower_figs/}{section_gpio_clock.tex}
\subimport{line_folower_figs/}{appendix_telemetry_video.tex}
\subimport{line_folower_figs/}{appendix_long_term_goal.tex}
\subimport{line_folower_figs/}{appendix_robot_motion.tex}
\subimport{line_folower_figs/}{appendix_photo3_calibration.tex}
\subimport{line_folower_figs/}{appendix_bench_history.tex}
\subimport{line_folower_figs/}{appendix_historical_cad.tex}
\end{document}
